% ============================================================
%  sec-1-usecase-dataset.tex  —  4.1 Use Case, Dataset, and
%  Project Structure
% ============================================================

\section{Use Case Definition and Dataset Construction}
\label{sec:usecase-dataset}


The migration target for this project is a concrete, industrially
grounded use case: a \term{Multilingual Technical Support Assistant}
for \acrshort{mena} Telecom Infrastructure. The system is designed
to assist enterprise clients of \CompanyName{} — telecom operators,
\acrshort{isp}s, and enterprise \acrshort{it} teams — in resolving
technical issues related to 5G infrastructure, Huawei network
equipment, optical transport networks, and enterprise \acrshort{it}
systems. Supported languages are English, Arabic, and French,
reflecting the linguistic profile of \CompanyName{}'s Northern Africa
customer base.

This use case was selected for three concrete reasons. First, Nemotron
3 Nano's native multilingual support (15+ languages) and 1M-token
context window align directly with the requirements of multi-document
technical queries and multi-turn support conversations. Second, the
3.6B active parameter count fits within the 64~GB \acrshort{hbm2}
of a single Ascend 910C, enabling single-node inference without
tensor parallelism. Third, the domain-specificity of technical support
provides clear, measurable evaluation criteria — keyword coverage,
command correctness, and language accuracy — rather than the
subjective metrics of general-purpose assistants.

The use case positions the fine-tuned model as the intelligence layer
of a production support platform, with the inference server exposing
an OpenAI-compatible \acrshort{rest} \acrshort{api} that can be
integrated into ticketing systems, chat interfaces, and
\acrshort{cli} tooling without modification.

\subsection{Dataset Blend}
\label{subsec:dataset}

The \acrshort{sft} dataset is constructed from four publicly
available HuggingFace datasets, blended by weight to reflect the
target language and domain distribution. \tabref{dataset-blend}
documents each source.

\begin{table}[htbp]
  \centering
  \caption{SFT dataset blend for the telecom support use case}
  \label{tab:dataset-blend}
  \begin{tabularx}{\textwidth}{P{3.5cm} L c P{2cm} r}
    \toprule
    \tableheader{Dataset & HuggingFace Path & Lang. & Size & Weight} \\
    \midrule
    Salesforce DialogStudio
      & \texttt{Salesforce/dialogstudio}
      & EN & $\sim$50K & 0.40 \\
    \rowcolor{TableRowOdd}
    Arabic Alpaca
      & \texttt{kunishou/ArabicAlpaca}
      & AR & $\sim$52K & 0.30 \\
    Open-Platypus
      & \texttt{garage-bAInd/Open-Platypus}
      & EN/FR & $\sim$25K & 0.20 \\
    \rowcolor{TableRowOdd}
    OPUS-100 (fr-ar)
      & \texttt{Helsinki-NLP/opus-100}
      & FR/AR & $\sim$10K & 0.10 \\
    \midrule
    \textbf{Total}
      & & & $\sim$137K & 1.00 \\
    \bottomrule
  \end{tabularx}
\end{table}

Each dataset undergoes a distinct formatting transformation
implemented in \filepath{data/prepare\_datasets.py}. The
\textbf{DialogStudio TechSupport} split contains multi-turn
technical support conversations in OpenAI dialog format; each
dialog is parsed by iterating over \texttt{roles} fields and
mapping \texttt{USR} tokens to \texttt{user} turns and \texttt{SYS}
tokens to \texttt{assistant} turns. The \textbf{Arabic Alpaca}
dataset uses an instruction-input-output structure; the
\texttt{instruction} and optional \texttt{input} fields are
concatenated into the \texttt{user} turn, and the \texttt{output}
becomes the \texttt{assistant} turn. The \textbf{Open-Platypus}
dataset is filtered to retain only samples containing technical
keywords (network, protocol, server, system, réseau, etc.) before
formatting. The \textbf{OPUS-100 fr-ar} translation pairs are
converted into synthetic bilingual QA samples where a French
technical sentence becomes the user query and the Arabic translation
becomes the assistant response, filtered to retain only
telecom-relevant vocabulary.

All samples are formatted to a unified OpenAI chat \acrshort{json}
schema with system, user, and assistant roles, as described in
Chapter 3.

Three language-specific system prompts are defined in
\filepath{data/prepare\_datasets.py} and injected based on
dataset source:

\begin{itemize}
  \item \texttt{SYSTEM\_PROMPT\_EN} — English prompt positioning the
        model as a Huawei telecom support specialist
  \item \texttt{SYSTEM\_PROMPT\_AR} — Arabic prompt with equivalent
        framing in \acrshort{msa}
  \item \texttt{SYSTEM\_PROMPT\_FR} — French prompt for
        French-language client interactions
\end{itemize}

During inference, the server automatically detects the language of
the incoming user query (via Unicode Arabic character ratio and
French lexical markers) and injects the appropriate system prompt,
ensuring consistent role framing without requiring the client
application to supply a system prompt.

\subsection{Dataset Preparation Pipeline}
\label{sec:data-prep-pipeline}

The dataset preparation executes in two stages. In the first stage, raw
datasets are downloaded and formatted to \texttt{.jsonl} using
\filepath{data/prepare\_datasets.py}.

The second stage runs the Nemotron SFT data pipeline on the formatted
files to apply the chat template and produce packed Parquet shards.
